
\documentclass{article}
\usepackage{colortbl}
\usepackage{makecell}
\usepackage{multirow}
\usepackage{supertabular}

\begin{document}

\newcounter{utterance}

\centering \large Interaction Transcript for game `ifeval', experiment `ifeval', episode 477 with reflexion{-}Qwen3.5{-}4B.
\vspace{24pt}

{ \footnotesize  \setcounter{utterance}{1}
\setlength{\tabcolsep}{0pt}
\begin{supertabular}{c@{$\;$}|p{.15\linewidth}@{}p{.15\linewidth}p{.15\linewidth}p{.15\linewidth}p{.15\linewidth}p{.15\linewidth}}
    \# & \multicolumn{2}{c}{Player} && \multicolumn{2}{c}{Game Master} \\
    \hline

    \theutterance \stepcounter{utterance}  
    & & & \multicolumn{4}{p{0.6\linewidth}}{
        \cellcolor[rgb]{0.9,0.9,0.9}{
            \makecell[{{p{\linewidth}}}]{
                \texttt{\tiny{[P1$\langle$GM]}}
                \texttt{You have attempted this same task before and did not succeed. These are the plans you wrote after those attempts.} \\
\texttt{Trial \#1: New plan: I failed because I used the letter "a" four times in the first stanza ("alone", "a", "dawn", "lane", "and", "pain", "plan"). I should have counted the letter "a" before writing each line to ensure the total count never exceeded one. Next time, I will write the entire song in my head or on scratchpad first, strictly limiting the single "a" to a specific word (like "a" or "Layton") and avoiding it in all other lines, then generate the full lyrics only after confirming the constraint is met.} \\
\\ 
\\ 
\texttt{Write a song about Layton, making sure to use the letter "a" at most once.} \\
            }
        }
    }
     \\ \\

    \theutterance \stepcounter{utterance}  
    & \multicolumn{4}{p{0.6\linewidth}}{
        \cellcolor[rgb]{0.561,0.737,0.561}{
            \makecell[{{p{\linewidth}}}]{
                \texttt{\tiny{[P1$\rangle$GM]}}
                \texttt{Layton walks the street,} \\
\texttt{No rain, no heat.} \\
\texttt{He sees a cat,} \\
\texttt{Then runs away.} \\
\texttt{Layton sleeps at night,} \\
\texttt{No light, no sight.} \\
\texttt{Just quiet air,} \\
\texttt{And empty chair.} \\
\texttt{Layton smiles today,} \\
\texttt{No gray, no way.} \\
\texttt{He finds a star,} \\
\texttt{And leaves it there.} \\
            }
        }
    }
    & & \\ \\

    \theutterance \stepcounter{utterance}  
    & & & \multicolumn{2}{p{0.3\linewidth}}{
        \cellcolor[rgb]{0.9,0.9,0.9}{
            \makecell[{{p{\linewidth}}}]{
                \texttt{\tiny{[GM$|$GM]}}
                {'keywords:letter_frequency': {'let_relation': 'less than', 'letter': 'a', 'let_frequency': 2}}
            }
        }
    }
    & & \\ \\

    \theutterance \stepcounter{utterance}  
    & & & \multicolumn{2}{p{0.3\linewidth}}{
        \cellcolor[rgb]{0.9,0.9,0.9}{
            \makecell[{{p{\linewidth}}}]{
                \texttt{\tiny{[GM$|$GM]}}
                \texttt{game\_result = LOSE} \\
            }
        }
    }
    & & \\ \\

    \theutterance \stepcounter{utterance}  
    & & \multicolumn{4}{p{0.6\linewidth}}{
        \cellcolor[rgb]{0.95,0.95,0.95}{
            \makecell[{{p{\linewidth}}}]{
                \texttt{\tiny{[P1$\leftrightarrow$P]}}
                \texttt{Post{-}episode Reflexion (not sent to GM):} \\
\texttt{New plan: I failed because I did not track the cumulative count of the letter "a" across the entire poem, only checking locally within stanzas. Next time, I will draft the entire song first, explicitly listing every word and counting the total occurrences of "a" before writing a single line. I will reserve the letter "a" for only one specific word (e.g., "Layton" or "a") and ensure no other word contains it, verifying the final count is exactly one or zero before submitting.} \\
            }
        }
    }
    & \\ \\

\end{supertabular}
}

\end{document}
